\documentclass{article}
\usepackage[default]{arsenal}
\PassOptionsToPackage{hyphens}{url}
\usepackage{hyperref}
\urlstyle{rm}
\usepackage{microtype}
\usepackage[ukrainian, english]{babel}
\usepackage{natbib}
\setcounter{secnumdepth}{0}
\makeatletter
\def\GetFileInfo#1{%
  \def\filename{#1}%
  \def\@tempb##1 ##2 ##3\relax##4\relax{%
    \def\filedate{##1}%
    \def\fileversion{##2}%
    \def\fileinfo{##3}}%
  \edef\@tempa{\csname ver@#1\endcsname}%
  \expandafter\@tempb\@tempa\relax? ? \relax\relax}
\makeatother
\GetFileInfo{arsenal.sty}
\newcommand{\abc}{abcdefghijklmnopqrstuvwxyz}
\newcommand{\ABC}{ABCDEFGHIJKLMNOPQRSTUVWXYZ}
\newcommand{\alphabeta}{\alpha\beta\gamma\delta\epsilon\varepsilon\zeta\eta\theta\vartheta\iota\kappa\varkappa\lambda\mu\nu\xi o\pi\varpi\rho\varrho\sigma\varsigma\tau\upsilon\phi\varphi\chi\psi\omega}
\newcommand{\AlphaBeta}{\Gamma\Delta\Theta\Lambda\Xi\Pi\Sigma\Upsilon\Phi\Psi\Omega}
\providecommand\pkg[1]{\textit{#1}}

\usepackage{unicodefonttable}

\newcommand\device{{\centering\textaldine\textaldine\textaldine\par}}


\begin{document}
\selectlanguage{english}
\title{Sample of Arsenal fonts}
\author{Boris Veytsman}
\date{Arsenal package version \fileversion, \filedate}
\maketitle
\tableofcontents


\clearpage

\section{Address by President Volodymyr Zelenskyy, September 2, 2023}



\textit{Dear Ukrainians, I wish you good health!}

This Saturday, three cities are celebrating their day. The cities that
are indispensable part of Ukraine.

\textsc{Odesa.} Our support in the south, on the Black Sea. A city that,
together with Ukraine, has and will always have global significance. A
port on which the lives of various nations depend---from Ukrainian
exports through the Great Odesa. A city of culture that knows how to
be interesting to everyone and respects everyone.

We have defended Odesa from destruction. Because the Russian regime is
incapable of bringing anything other than degradation. And we will
return security to Odesa. Odesa has always been a place where you feel
lightness and happiness. Odesa will remain this way. Congratulations
on your Day!

\textsc{Sumy.} Our outpost in the northeast. Every year, on the second Saturday
of September, it celebrates its day. And it will always celebrate it
as a free, Ukrainian city.

During these times---the times of war---unfortunately, we often
receive reports of Russian terror from Sumy. About shelling, missiles,
and bombs. About constant attempts by Russian sabotage groups to
infiltrate the region.

But despite everything, Sumy region is alive. Sumy is alive and gives
strength to the entire region. And when I was in the city, I felt that
there is faith there. Faith that evil will not prevail. Faith in
people. Faith in Ukraine. Faith that we will definitely get through
this time. And we will win. It will be so. Sumy, congratulations!

And the third city---\textsc{Lysychansk.} A city that Ukraine still needs to
reclaim along with the entire Luhansk region.

Today, no one can specify a date when the city will be free again. But
everyone who fights and works for Ukraine is doing everything possible
so that our cities and villages currently under occupation can once
again experience normal and free life.

Lysychansk has always been one of the pillars of the east of our
state, one of the key cities. A proud city! It will remain so. A city
that knows how to work and is rightfully proud of its
achievements. Together with Ukraine, it's all possible. With our
strength, unity, and our ability to take care of each other – all
cities together, all villages, all people.

And one more thing worth mentioning.

Undoubtedly, we will defend Ukraine and restore freedom to all our
land. Each of us feels that this will be a Ukraine with different
rules. The borders are the same. Democracy is probably just as
turbulent. Freedom is one of the greatest in Europe, as always.

But without a doubt, there will be no more decades-long ``business as
usual'' for those who plundered Ukraine and put themselves above the
law and any rules. And I thank the Ukrainian law enforcement for their
determination to bring every case stalled for decades to a just
conclusion. The law must work. It is so. It will be so.

\textit{Glory to Ukraine!}

\device

\bigskip
{\footnotesize From
  \url{https://www.president.gov.ua/en/news/ciyeyi-suboti-svij-den-vidznachayut-tri-mista-bez-yakih-ne-u-85345},
  vi\-sited on September 3, 2023.} 


\clearpage

\selectlanguage{ukrainian}

\section{Звернення Президента Володимира Зеленського, 2 вересня 2023
  року}

\textit{Бажаю здоров’я, шановні українці, українки!}

Цієї суботи три міста відзначають свій день. Міста, без яких не уявити Україну.

\textsc{Одеса.} Наша опора на півдні, на Чорному морі. Місто, яке разом з
Україною має та завжди матиме глобальне значення. Порт, від якого
залежить життя різних народів "--- від українського експорту через Велику
Одесу. Місто культури, яке вміє бути цікавим для всіх та вміє поважати
кожного й кожну. 

Ми захистили Одесу від знищення. Бо нічого іншого, крім деградації,
російський режим нездатний принести. І ми повернемо Одесі
безпеку. Завжди Одеса була такою, що в ній відчуваєш легкість і
щастя. Такою Одеса і буде. Вітаю вас із вашим Днем!

\textsc{Суми.} Наш форпост на північному сході. Щороку в другу суботу вересня
відзначає свій день. І буде відзначати завжди вільним, завжди
українським.

У цей час "--- час війни "--- із Сумщини, на жаль, дуже часто надходять
повідомлення про російський терор. Про обстріли, ракети, бомби. Про
постійні намагання російських ДРГ зайти в область.

Але, попри все, Сумщина живе. Суми живуть, дають силу всій області. І
коли я був у місті, я відчув, що там є віра. Віра, що зло не стане
вищим. Віра в людей. Віра в Україну. Віра в те, що ми обовʼязково
пройдемо цей час. І переможемо. Так буде. Суми, вітаю вас!

І третє місто "--- \textsc{Лисичанськ.} Місто, яке Україні ще належить
повернути разом з усією Луганщиною.

Сьогодні ніхто не назве конкретної дати, коли місто буде знову
вільним. Але кожен, хто воює та працює заради України, робить усе
можливе, щоб наші міста й села, які зараз в окупації, могли знову
відчути нормальне та вільне життя.

Лисичанськ завжди був однією з опор сходу нашої держави, одним із
ключових міст. Горде місто! Таким і буде. Місто, яке вміє працювати та
по праву пишається своїми результатами. Разом з Україною це все
дається. З нашою силою, нашою єдністю та нашим умінням дбати одне про
одного "--- усі міста разом, усі села, усі люди.

І ще одне, про що варто сказати.

Безумовно, ми захистимо Україну та повернемо свободу всій землі. І
кожен із нас відчує, що це буде Україна інших правил. Кордони "--- ті
самі. Демократія "--- напевно, така ж бурхлива. Свобода "--- одна з
найбільших у Європі, як і завжди.

Але точно без багатолітнього «як завжди» щодо тих, хто грабував
Україну та ставив себе вище, ніж закон і будь-які правила. І я дякую
українським правоохоронцям за рішучість довести до справедливого
результату кожну "--- кожну зі справ, які десятиліттями
гальмувались. Закон повинен працювати. Так є. Так буде.

\textit{Слава Україні!}

\device


\bigskip
{\footnotesize З
  \url{https://www.president.gov.ua/ua/news/ciyeyi-suboti-svij-den-vidznachayut-tri-mista-bez-yakih-ne-u-85345},
  відвідано 3 вересня 2023 року.\par} 

\clearpage

\selectlanguage{english}
\section{Math from \citep{Hartke06,
  free-math-font-survey}}

\textbf{Theorem 1 (Residue Theorem).}
Let $f$ be analytic in the region $G$ except for the isolated singularities $a_1,a_2,\ldots,a_m$. If $\gamma$ is a closed rectifiable curve in $G$ which does not pass through any of the points $a_k$ and if $\gamma\approx 0$ in $G$ then
\[
\frac{1}{2\pi i}\int_\gamma f = \sum_{k=1}^m n(\gamma;a_k) \text{Res}(f;a_k).
\]

\textbf{Theorem 2 (Maximum Modulus).}
\emph{Let $G$ be a bounded open set in $\mathbb{C}$ and suppose that $f$ is a continuous function on $G^-$ which is analytic in $G$. Then}
\[
\max\{|f(z)|:z\in G^-\}=\max \{|f(z)|:z\in \partial G \}.
\]

\selectlanguage{ukrainian}

\section{Математика з \citep{Hartke06, free-math-font-survey}}


\textbf{Теорема 1 (Теорема про залишки).}
Нехай $f$ аналітична в області $G$ за винятком ізольованих
сингулярностей $a_1,a_2,\ldots,a_m$. Якщо $\gamma$ є замкнута крива  в $G$, що
може бути спрямована, яка не проходить скрізь жодну з точок
$a_k$,  і якщо $\gamma\approx 0$ в $G$, то
\[
\frac{1}{2\pi i}\int_\gamma f = \sum_{k=1}^m n(\gamma;a_k) \text{Res}(f;a_k).
\]

\textbf{Теорема 2 (Максимальне значення).}
\emph{Нехай $G$ є обмежена множина в $\mathbb{C}$, і нехай $f$ є
  безперервна функція на $G^-$, аналітична в $G$. Тоді}
\[
\max\{|f(z)|:z\in G^-\}=\max \{|f(z)|:z\in \partial G \}.
\]

\selectlanguage{english}

\section{Alphabets}

\bgroup
\setlength{\parindent}{0pt}
\setlength{\parskip}{1ex}

Uppercase and math\\
\ABC\quad \textit{\ABC} \quad $\ABC$

Lowercase and math\\
\abc\quad\textit{\abc} \quad $\abc$ \quad 0123456789\quad $01234567890$

Greek\\
$\AlphaBeta$ \quad $\alphabeta$ \quad $\ell\wp\aleph\infty\propto\emptyset\nabla\partial\mho\imath\jmath\hslash\eth$

Lowercase Greek and math\\
$\abc\quad \alphabeta$

Uppercase Greek and math\\
$\ABC\quad \AlphaBeta$

Greek and misc\\
$\mathrm{A} \Lambda \Delta \nabla \mathrm{B C D} \Sigma \mathrm{E F} \Gamma \mathrm{G H I J K L M N O} \Theta \Omega \mho \mathrm{P} \Phi \Pi \Xi \mathrm{Q R S T U V W X Y} \Upsilon \Psi \mathrm{Z} $  $ \quad 1234567890 $


Mathbold\\
\textbf{\ABC}\quad $\mathbf{\ABC}$\\
\textbf{\abc}\quad $\mathbf{\abc}$

Math and symbols\\
$a\alpha b \beta c \partial d \delta e \epsilon \varepsilon f \zeta \xi g \gamma h \hbar \hslash \iota i \imath j \jmath k \kappa \varkappa l \ell \lambda m n \eta \theta \vartheta o \sigma \varsigma \phi \varphi \wp p \rho \varrho q r s t \tau \pi u \mu \nu v \upsilon w \omega \varpi x \chi y \psi z$ \linebreak[3] $\infty \propto \emptyset \varnothing \mathrm{d}\eth \backepsilon$

Mathcal\\
$\mathcal{\ABC}$

Mathbb\\
$\mathbb{\ABC}$

Mathscr\\
$\mathscr{\ABC}$

Uppercase mathfrak\\
$\mathfrak{\ABC}$

Lowercase mathfrak\\
$\mathfrak{\abc}$

Bold math\\
{\boldmath $\alpha + b = 27$}

Primes:
$f', f'', f'''$.
\egroup




\section{Font table}


\displayfonttable{Arsenal-Regular.otf}

\bibliography{arsenal}
\bibliographystyle{plainnat}

\end{document}
